\chapter{Wprowadzenie do teorii Seiberga--Wittena i zastosowań w wymiarze cztery}
\label{ch:tg-sw}

Rachunek grafów plumbingowych daje nam formy przecięcia,
grupy homologii i sygnatury. Nie odpowiada jednak na pytanie,
czy dwie homeomorficzne czterorozmaitości są dyfeomorficzne.
Do rozstrzygania tego typu problemów wykorzystuje się
niezmienniki teorii cechowania. Naszym celem jest tu
precyzyjne wprowadzenie niezbędnych obiektów, wyprowadzenie
rachunków topologicznych i wyjaśnienie zastosowania tych
niezmienników do wymiernej kontrakcji. Zaawansowane
fakty analityczne jawnie przyjmujemy z literatury.

\section{Formy samodualne i klasy charakterystyczne}
\label{sec:tg-sw-selfdual}

Niech $X$ będzie zamkniętą, spójną, zorientowaną,
gładką czterorozmaitością z metryką riemannowską $g$.
Operator Hodge'a $*_g:\Omega^2(X)\to\Omega^2(X)$ spełnia
$*_g^2=\mathrm{id}$. Rozkłada zatem wiązkę form na podwiązki
$\Lambda^2=\Lambda^2_+\oplus\Lambda^2_-$, będące wiązkami
własnymi dla wartości własnych odpowiednio $+1$ i $-1$.
Dla formy stopnia dwa $\eta$ oznaczamy
$\eta^+:=(\eta+*_g\eta)/2$.

\begin{stwierdzenie}[Formy samodualne i liczba $b_2^+$]
Niech $X$ będzie zamkniętą, spójną, zorientowaną
czterorozmaitością riemannowską. Przestrzeń harmonicznych
form samodualnych ma wymiar $b_2^+(X)$, równy dodatniemu
indeksowi bezwładności formy przecięcia na $H^2(X;\mathbb R)$.
Analogicznie wymiar harmonicznych form antysamodualnych
jest równy $b_2^-(X)$.
\end{stwierdzenie}

\begin{proof}
Z twierdzenia Hodge'a każda klasa de Rhama stopnia dwa
ma dokładnie jednego reprezentanta harmonicznego.
Na przestrzeni harmonicznych form operator $*_g$ zachowuje
harmoniczność, bo komutuje z laplasjanem Hodge'a.
Dla harmonicznej formy samodualnej $\omega$ zachodzi
$\int_X\omega\wedge\omega=\int_X\omega\wedge*_g\omega
=\int_X|\omega|^2\,\mathrm d\mu_g>0$, jeśli $\omega\ne0$.
Dla formy antysamodualnej analogiczna całka jest ujemna.
Dla harmonicznych form $\omega_+$ i $\omega_-$
z odpowiednich podprzestrzeni mamy punktowo
$\langle\omega_+,\omega_-\rangle=0$,
ponieważ operator Hodge'a jest samosprzężony
i jego wartości własne są różne.
Stąd $\omega_+\wedge\omega_-=
-\omega_+\wedge *_g\omega_-=0$.
Obie przestrzenie są więc ortogonalne względem
formy przecięcia, a ich suma daje całe $H^2(X;\mathbb R)$.
Tak więc ich wymiary są właśnie dodatnim i ujemnym
indeksem tej formy.
\end{proof}

\begin{definicja}[Klasa charakterystyczna]
Niech $X$ będzie jak wyżej. Klasę $c\in H^2(X;\mathbb Z)$
nazywamy \emph{charakterystyczną}, jeżeli
$\langle c\smile x,[X]\rangle\equiv
\langle x\smile x,[X]\rangle\pmod2$
dla każdego $x\in H^2(X;\mathbb Z)$.
Jeżeli nie występuje torsja w $H^2(X;\mathbb Z)$,
warunek ten jest równoważny równości
$c\bmod2=w_2(TX)\in H^2(X;\mathbb Z/2)$,
na mocy wzoru Wu. Istotnie, brak torsji w $H^2(X;\mathbb Z)$
pociąga za sobą brak torsji w $H_1(X;\mathbb Z)$
z uniwersalnego twierdzenia o współczynnikach.
Z dualności Poincarégo także $H^3(X;\mathbb Z)$
nie ma torsji. Długi ciąg współczynników dla
$0\to\mathbb Z\xrightarrow{\cdot2}\mathbb Z
\to\mathbb Z/2\to0$ pokazuje więc, że redukcja
$H^2(X;\mathbb Z)\to H^2(X;\mathbb Z/2)$
jest surjektywna. Wzór Wu i nieosobliwość
parowania Poincarégo modulo dwa dają żądaną
równoważność. W obecności torsji sama zgodność
parowań z klasami całkowitymi nie musi wykrywać całej
informacji modulo dwa. Dlatego dla wiązki wyznacznikowej
struktury $\mathrm{Spin}^c$ wymagamy jawnie
$c_1(L)\bmod2=w_2(TX)$.
\end{definicja}

\begin{przyklad}[Klasy charakterystyczne na sumie spójnej]
Na $X=\mathbb{CP}^2\#k\overline{\mathbb{CP}}^{\,2}$
przyjmijmy bazę $H,E_1,\ldots,E_k$ o kwadratach
$H^2=1$, $E_i^2=-1$ i zerowych przecięciach mieszanych.
Ponieważ macierz przecięcia jest diagonalna z nieparzystymi
wpisami, klasa $c=aH-\sum b_iE_i$ jest charakterystyczna
wtedy i tylko wtedy, gdy wszystkie liczby $a,b_i$ są
nieparzyste. W szczególności warunek charakterystyczności
ogranicza zbiór klas, które mogą być podstawą
niezmienników Seiberga--Wittena, ale nie orzeka,
że każda taka klasa daje niezerowy niezmiennik.
\end{przyklad}

\section{Struktura \texorpdfstring{$\mathrm{Spin}^c$}{Spinc} i równania}
\label{sec:tg-sw-equations}

\begin{definicja}[Dane struktury $\mathrm{Spin}^c$]
Niech $X$ będzie zamkniętą zorientowaną gładką
czterorozmaitością. Struktura $\mathrm{Spin}^c$ na $X$
(po wyborze metryki) jest podniesieniem głównej wiązki
dodatnio zorientowanych ortonormalnych układów
odniesienia do grupy
\[
\mathrm{Spin}^c(4)=
(\mathrm{Spin}(4)\times S^1)/\langle(-1,-1)\rangle.
\]
Podniesienie oznacza główną wiązkę tej grupy
z odwzorowaniem do wiązki układów odniesienia
zgodnym z homomorfizmem $[a,z]\mapsto\pi(a)$,
gdzie $\pi:\mathrm{Spin}(4)\to SO(4)$ jest
podwójnym nakryciem. Reprezentacje półspinorowe
$[a,z]\mapsto z\rho_\pm(a)$ indukują zespolone
wiązki spinorowe $S^\pm$ rangi dwa. Reprezentacja
wyznacznikowa $[a,z]\mapsto z^2$ daje wiązkę
liniową $L=\det S^+\cong\det S^-$.
Klasa $c_1(L)$ spełnia $c_1(L)\bmod2=w_2(TX)$.
Przy ustalonej metryce i strukturze $\mathrm{Spin}^c$
połączenie unitarne $A$ na $L$ wyznacza operator Diraca
$D_A:\Gamma(S^+)\to\Gamma(S^-)$, zadaną lokalnie
w ortonormalnej bazie $(e_j)_{j=1}^4$ wzorem
$D_A\psi=\sum_{j=1}^4\rho(e^j)\nabla^A_{e_j}\psi$,
gdzie $(e^j)$ jest bazą dualną, $\rho$ oznacza
mnożenie Clifforda, a $\nabla^A$ jest połączeniem
spinorowym indukowanym przez $A$ i połączenie
Levi-Civity.
W szczególności zmiana $A$ o formę
$a\in i\Omega^1(X)$ zmienia połączenie spinorowe
o $a/2$, gdyż na spinorach środkowy okrąg
działa z wagą jeden, a na wyznaczniku z wagą dwa.
\end{definicja}

\begin{uwaga}
Konstrukcja wiązek spinorowych wymaga algebry Clifforda
i pojęcia struktur spinorowych. Odnotujmy standardową
własność klasyfikacyjną: po ustaleniu jednej struktury
$\mathfrak s$ pozostałe otrzymuje się jednoznacznie
(z dokładnością do izomorfizmu) przez skręcenie
o zespoloną wiązkę liniową $M$. Wtedy
$S^\pm_{\mathfrak s\otimes M}=S^\pm_{\mathfrak s}\otimes M$,
a z rangi dwa wynika
$L_{\mathfrak s\otimes M}=L_{\mathfrak s}\otimes M^{\otimes2}$.
Zatem $c_1(L_{\mathfrak s\otimes M})=c_1(L_{\mathfrak s})+2c_1(M)$.
Jeśli istnieje niezerowa klasa $x$ z $2x=0$,
dwie różne struktury mogą mieć tę samą klasę
wyznacznikową. Samą klasyfikację podniesień
przyjmujemy tutaj jako standardowy wynik teorii
struktur $\mathrm{Spin}^c$.
W dalszym ciągu korzystamy z podstawowych własności
operatora Diraca oraz z twierdzenia indeksowego
Atiyaha--Singera; nie odtwarzamy tu wszystkich narzędzi
analitycznych teorii cechowania.
\end{uwaga}

\begin{definicja}[Równania Seiberga--Wittena]
\label{def:tg-sw-eqn}
Niech $g$ będzie metryką riemannowską na zamkniętej,
spójnej, zorientowanej czterorozmaitości $X$, a
$\mathfrak s$ strukturą $\mathrm{Spin}^c$ o wiązce
wyznacznikowej $L$. Niech $\eta\in\Omega^2_+(X)$
będzie ustaloną rzeczywistą formą samodualną.
Szukamy pary $(A,\psi)$, gdzie $A$ jest połączeniem
unitarnym na $L$, a $\psi\in\Gamma(S^+)$ spinorem,
spełniającej
\[
D_A\psi=0,\qquad F_A^+=q(\psi)+i\eta.
\]
Tutaj $F_A\in i\Omega^2(X)$ jest krzywizną połączenia,
$F_A^+$ jej częścią samodualną, a
$q:S^+\to i\Lambda^2_+$ odwzorowaniem kwadratowym
określonym, w konwencji definicji
\ref{def:tg-sw12-q}, przez
\[
\rho(q(\psi))=\psi\otimes\psi^*
-\tfrac12|\psi|^2\operatorname{id}_{S^+}.
\]
Tutaj $\rho:i\Lambda^2_+\to\operatorname{Herm}_0(S^+)$
jest izomorfizmem rzeczywistych wiązek indukowanym
przez mnożenie Clifforda. Iloczyn hermitowski
jest liniowy w pierwszym argumencie, więc
$(\psi\otimes\psi^*)(v)=\langle v,\psi\rangle\psi$.
Przy tej normalizacji człon krzywizny we wzorze
Weitzenböcka ma postać $\frac12\rho(F_A^+)$.
Szczegółowe rachunki algebraiczne i oszacowania
podajemy w rozdziale \ref{ch:tg-sw-analiza-uzup}.
\end{definicja}

\begin{definicja}[Przestrzeń moduli]
\emph{Grupą cechowania} nazywamy
$\mathcal G=C^\infty(X,S^1)$.
Działa ona na parach $(A,\psi)$ przez
$u\cdot(A,\psi)=(A-2u^{-1}\mathrm du,u\psi)$,
przy standardowej konwencji połączenia wyznacznikowego.
\emph{Przestrzenią moduli} $\mathcal M_{\mathfrak s}(g,\eta)$
nazywamy zbiór rozwiązań równań z definicji
\ref{def:tg-sw-eqn}, podzielony przez działanie $\mathcal G$.
\end{definicja}

Działanie jest dobrze określone na zbiorze rozwiązań:
jeśli $A'=A-2u^{-1}\mathrm du$, to
$\nabla^{A'}(u\psi)=u\nabla^A\psi$,
ponieważ zmiana połączenia spinorowego wynosi
$-u^{-1}\mathrm du$. Zatem $D_{A'}(u\psi)=uD_A\psi$.
Forma $u^{-1}\mathrm du$ jest lokalnie równa
$i\,\mathrm df$ dla rzeczywistego $f$,
więc jest domknięta i $F_{A'}=F_A$.
Ponadto $q(u\psi)=q(\psi)$, gdyż $|u|=1$.
Oba równania pozostają zatem spełnione.

\begin{uwaga}[Rozwiązania redukowalne]
Pary z $\psi=0$ nazywamy \emph{redukowalnymi}.
Na spójnej rozmaitości stabilizator konfiguracji
redukuje się do stałych funkcji z $S^1$:
warunek zachowania $A$ wymusza $\mathrm du=0$,
a gdy $\psi$ nie znika identycznie, warunek
$u\psi=\psi$ wymusza dodatkowo $u=1$.
Rozwiązania redukowalne spełniają
równanie $F_A^+=i\eta$.
Dla $b_2^+>0$ można generycznie dobrać perturbację
tak, aby uniknąć rozwiązań redukowalnych. Przy odpowiedniej
transwersalności przestrzeń moduli jest wtedy gładka,
a twierdzenie o zwartości zapewnia jej zwartość.
Przy $b_2^+>1$ generyczna jednoparametrowa rodzina danych
może dodatkowo ominąć ścianę redukowalnych. Nie jest to jednak wyłącznie argument
kombinatoryczny: wymaga analizy eliptycznej, transwersalności
i twierdzenia o zwartości.
\end{uwaga}

\section{Wymiar oczekiwany i definicja niezmiennika}
\label{sec:tg-sw-dimension}

\begin{twierdzenie}[Oczekiwany wymiar przestrzeni moduli]
\label{thm:tg-sw-dimension}
Dla zamkniętej, spójnej, zorientowanej,
gładkiej czterorozmaitości $X$ i struktury
$\mathrm{Spin}^c$ z klasą $c=c_1(L)$ rzeczywisty
wymiar oczekiwany przestrzeni moduli wynosi
\[
d(\mathfrak s)=\frac{\langle c^2,[X]\rangle-2e(X)-3\sigma(X)}4.
\]
Jest on liczbą całkowitą. Dla generycznych danych bez rozwiązań
redukowalnych, przy transwersalności, jest to
wymiar niepustej przestrzeni moduli.
Ujemny indeks oznacza wtedy pustą przestrzeń.
Warunek $b_2^+(X)>1$ potrzebny w dalszej części
dotyczy niezmienniczości, a nie samego rachunku indeksu.
\end{twierdzenie}

\begin{proof}[Rachunek indeksu; użycie twierdzenia Atiyaha--Singera]
Kompleks deformacyjny teorii cechowania zawiera część
Diraca oraz część rzeczywistą z operatorem
$\mathrm d^*\oplus\mathrm d^+:\Omega^1(X)\to
\Omega^0(X)\oplus\Omega^2_+(X)$.
Twierdzenie Atiyaha--Singera daje zespolony indeks
operatora Diraca jako ewaluację części stopnia cztery
klasy $e^{c/2}\widehat A(TX)$ na $[X]$.
Ponieważ $e^{c/2}=1+c/2+c^2/8$ i
$\widehat A(TX)=1-p_1(TX)/24$ w stopniach
nieprzekraczających czterech, a twierdzenie Hirzebrucha
daje $\langle p_1(TX),[X]\rangle=3\sigma(X)$,
otrzymujemy
$\operatorname{ind}_{\mathbb C}D_A
=(\langle c^2,[X]\rangle-\sigma(X))/8$.
Część rzeczywista wnosi do indeksu deformacji
$b_1(X)-b_0(X)-b_2^+(X)=b_1(X)-1-b_2^+(X)$.
Istotnie, jądro $\mathrm d^*\oplus\mathrm d^+$
składa się z harmonicznych 1-form: dla
$\mathrm d^+a=0$ dokładna 2-forma $\mathrm da$
jest antysamodualna, a z twierdzenia Stokesa
$0=\int_X\mathrm da\wedge\mathrm da
=-\|\mathrm da\|_{L^2}^2$. Stąd $\mathrm da=0$.
Jądro operatora sprzężonego składa się z par
$(f,\omega)$ z $\mathrm df+\mathrm d^*\omega=0$,
gdzie $\omega$ jest samodualna. Składniki są
ortogonalne w $L^2$, gdyż $\mathrm d^2 f=0$.
Zatem $f$ jest stała, a $\mathrm d^*\omega=0$
i samodualność wymuszają harmoniczność $\omega$.
Eliptyczność i niezmienniczość indeksu przy
zaburzeniach rzędu zero uzasadniamy w rozdziałach
\ref{ch:tg-sw-analiza} i \ref{ch:tg-sw-analiza-uzup}.
Łączny indeks rzeczywisty wynosi wobec tego
\[
2\operatorname{ind}_{\mathbb C}D_A+b_1-1-b_2^+
=\frac{\langle c^2,[X]\rangle-\sigma(X)}4+b_1-1-b_2^+.
\]
Ponieważ $e(X)=2-2b_1+b_2^++b_2^-$ i
$\sigma(X)=b_2^+-b_2^-$, prawa strona po uproszczeniu
jest równa wyrażeniu w tezie.
Całkowitość wynika z indeksowej interpretacji.
Nie dowodzimy tu samego twierdzenia Atiyaha--Singera ani
analitycznej regularności przestrzeni moduli.
\end{proof}

\begin{wniosek}[Parzystość wymiaru oczekiwanego]
\label{cor:tg-sw-index-parity}
Niech $X$ będzie zamkniętą, spójną, zorientowaną,
gładką czterorozmaitością, a $\mathfrak s$
ustaloną strukturą $\mathrm{Spin}^c$.
Wówczas
\[
d(\mathfrak s)\equiv
b_1(X)-1-b_2^+(X)\pmod2.
\]
W szczególności, gdy $b_1(X)=0$ i $b_2^+(X)$
jest parzyste, żadna struktura $\mathrm{Spin}^c$
nie ma oczekiwanego wymiaru zero.
\end{wniosek}

\begin{proof}
Niech $m=\operatorname{ind}_{\mathbb C}D_A\in\mathbb Z$.
Z rachunku indeksowego w dowodzie
Twierdzenia~\ref{thm:tg-sw-dimension}
otrzymujemy
$d(\mathfrak s)=2m+b_1(X)-1-b_2^+(X)$.
Po redukcji modulo dwa wyraz $2m$ znika,
co daje tezę. Gdy $b_1=0$ i $b_2^+$ jest
parzyste, prawa strona jest nieparzysta,
a więc $d(\mathfrak s)\ne0$.
Nie wnioskujemy stąd, że wszystkie
przestrzenie moduli są puste:
dla dodatnich nieparzystych wymiarów
teoria z dodatkowymi klasami kohomologii
wymaga osobnego sformułowania.
\end{proof}

\begin{przyklad}[Indeks dla $\mathbb{CP}^2\#\mathbb{CP}^2$]
\label{ex:tg-sw-index-parity}
Niech $X=\mathbb{CP}^2\#\mathbb{CP}^2$
z orientacją sumy spójnej.
Mamy $b_1=0$, $b_2^+=2$, $e(X)=4$
i $\sigma(X)=2$. Dla charakterystycznej
klasy $c=H_1+H_2$, gdzie
$H_1^2=H_2^2=1$ i $H_1H_2=0$,
dostajemy $c^2=2$, zatem
\[
d(\mathfrak s)=\frac{2-2\cdot4-3\cdot2}{4}
=-3.
\]
Dla dowolnej innej struktury
$\mathrm{Spin}^c$ na $X$ wymiar
oczekiwany pozostaje nieparzysty,
choć jego wartość może być inna.
Przykład ilustruje ograniczenie
parzystości, które jest silniejsze
od samego stwierdzenia
całkowitości indeksu.
\end{przyklad}

\begin{definicja}[Niezmiennik Seiberga--Wittena w wymiarze zero]
Niech $X$ będzie zamkniętą, spójną, zorientowaną,
gładką czterorozmaitością z $b_2^+(X)>1$.
Ustalmy orientację homologiczną, tj. orientację
$H^1(X;\mathbb R)\oplus H^2_+(X;\mathbb R)$,
gdzie drugi składnik jest dodatnio określoną
maksymalną podprzestrzenią formy przecięcia.
Orientację przenosimy na przestrzeń harmonicznych
form samodualnych przy zmianie metryki.
Dla $b_1(X)=0$ pozostaje tylko orientacja
drugiego składnika.
Dla struktury $\mathrm{Spin}^c$ o $d(\mathfrak s)=0$
i generycznych danych $(g,\eta)$ przestrzeń moduli
składa się ze skończonej liczby punktów, liczonych
z odpowiednimi znakami orientacji. Ich algebraiczną
liczbę oznaczamy przez $SW_X(\mathfrak s)\in\mathbb Z$.
Gdy $d(\mathfrak s)<0$, przyjmujemy $SW_X(\mathfrak s)=0$.
W przypadku $d(\mathfrak s)>0$ do zdefiniowania
liczbowych niezmienników używa się dodatkowych klas
kohomologii na przestrzeni moduli; nie są one częścią
powyższej definicji.
\end{definicja}

\begin{twierdzenie}[Niezmienniczość; wynik analityczny]
Dla $X$ jak w poprzedniej definicji, po ustaleniu
orientacji homologicznej, $SW_X(\mathfrak s)$ nie zależy
od wyboru danych generycznych $(g,\eta)$ i jest
niezmiennikiem gładkiej struktury na $X$, z dokładnością
do naturalnego działania dyfeomorfizmów na strukturach
$\mathrm{Spin}^c$ i na orientacji homologicznej.
\end{twierdzenie}

\begin{uwaga}[Przypadek $b_2^+=1$]
Dla $b_2^+(X)=1$ niezmiennik może zależeć od wyboru
\emph{komory}: podczas zmiany metryki i perturbacji
można przejść przez ścianę, na której występują
rozwiązania redukowalne. W związku z tym porównanie
z rozmaitościami wymiernymi, takimi jak
$\mathbb{CP}^2\#7\overline{\mathbb{CP}}^{\,2}$,
wymaga dodatkowej formuły zmiany komory i świadomego
wyboru tej komory. Nie wystarcza samo wyliczenie
$d(\mathfrak s)$.
\end{uwaga}

\section{Klasy podstawowe i ograniczenia geometryczne}
\label{sec:tg-sw-basic}

\begin{definicja}[Klasa podstawowa]
Dla $X$ z $b_2^+>1$ i dla wersji
niezmienników o wymiarze oczekiwanym zero klasę
charakterystyczną $c\in H^2(X;\mathbb Z)$
nazywamy \emph{klasą podstawową Seiberga--Wittena},
jeżeli istnieje struktura $\mathrm{Spin}^c$ o
$c_1(\det S^+)=c$ i $SW_X(\mathfrak s)\ne0$.
Przy obecności torsji odróżniamy klasę $c$ od samej
struktury $\mathrm{Spin}^c$.
\end{definicja}

\begin{twierdzenie}[Zerowanie przy dodatniej krzywiźnie skalarnej]
\label{thm:tg-sw-psc}
Jeżeli zamknięta, spójna, zorientowana,
gładka czterorozmaitość $X$ ma $b_2^+(X)>1$
i dopuszcza metrykę riemannowską o wszędzie dodatniej
krzywiźnie skalarnej, to wszystkie jej niezmienniki
Seiberga--Wittena z powyższej definicji są zerowe.
\end{twierdzenie}

\begin{proof}[Argument Weitzenböcka; krok analityczny wskazany jawnie]
Przy normalizacji z definicji \ref{def:tg-sw12-q}
tożsamość Weitzenböcka i równania dają po
scałkowaniu
\[
0=\int_X\left(
|\nabla_A\psi|^2+\tfrac{s_g}{4}|\psi|^2
+\tfrac14|\psi|^4+
\tfrac12\langle\rho(i\eta)\psi,\psi\rangle
\right)\,\mathrm d\mu_g.
\]
Ostatni wyraz jest rzeczywisty, ponieważ
$\rho(i\eta)$ jest hermitowski. Niech
$s_0=\min_X s_g>0$, a $C_\rho>0$ będzie
stałą z twierdzenia \ref{thm:tg-sw12-spinor-bound}.
Jeśli $2C_\rho\|\eta\|_{L^\infty}<s_0$, to
suma wyrazów kwadratowych w $\psi$ jest
ograniczona z dołu przez
$(s_0/4-C_\rho\|\eta\|_{L^\infty}/2)|\psi|^2$,
z dodatnim współczynnikiem. Wszystkie pozostałe
składniki są nieujemne, więc $\psi=0$.
Ściana perturbacji dopuszczających rozwiązania
redukowalne ma kowymiar $b_2^+(X)$
(stwierdzenie \ref{prop:tg-sw-wall}).
Wybieramy generyczną perturbację poza tą
ścianą w powyższej otwartej kuli.
Przestrzeń moduli jest wtedy pusta. Z twierdzenia
o niezmienniczości wynika zerowanie niezmiennika.
\end{proof}

\begin{twierdzenie}[Taubes; niezerowanie dla rozmaitości symplektycznych]
\label{thm:tg-sw-taubes}
Jeżeli $(X,\omega)$ jest zamkniętą, spójną,
symplektyczną czterorozmaitością z $b_2^+(X)>1$,
zorientowaną przez $\omega\wedge\omega$, to dla
kanonicznej struktury $\mathrm{Spin}^c$ indukowanej
przez zgodną strukturę prawie zespoloną $J$
oraz dla jej struktury sprzężonej niezmienniki
Seiberga--Wittena mają wartości bezwzględne jeden.
Jeżeli $K_\omega=-c_1(TX,J)$, klasy wyznacznikowe
tych struktur wynoszą odpowiednio $-K_\omega$
i $K_\omega$. Wynik przyjmujemy bez dowodu.
\end{twierdzenie}

Wynik o strukturze kanonicznej pochodzi z pracy
C.~H. Taubesa, \emph{The Seiberg--Witten invariants
and symplectic forms}, Math. Res. Lett. 1
(1994), 809--822, twierdzenie główne na s.~809.
Dla struktury prawie zespolonej wiązka
$S^+$ ma postać $\Lambda^{0,0}\oplus\Lambda^{0,2}$,
a jej wyznacznik jest izomorficzny z odwrotnością
wiązki kanonicznej. Stąd $c_1(L)=-K_\omega$.
Sprzężenie struktur $\mathrm{Spin}^c$ zmienia
znak klasy wyznacznikowej i zachowuje wartość
bezwzględną niezmiennika (symetria sprzężenia,
przyjmowana tutaj z teorii Seiberga--Wittena).
Obie struktury mają wymiar oczekiwany zero,
gdyż $K_\omega^2=c_1(TX,J)^2=2e(X)+3\sigma(X)$
na mocy wzoru dla rozmaitości prawie zespolonych
z twierdzenia \ref{thm:tg-rb-invariants}.
Przy torsji nie zastępujemy tych dwóch konkretnych
struktur dowolnymi strukturami o tych samych
klasach wyznacznikowych.

\begin{wniosek}
Zamknięta symplektyczna czterorozmaitość
$X$ z $b_2^+(X)>1$ nie dopuszcza metryki o wszędzie
dodatniej krzywiźnie skalarnej.
\end{wniosek}

\begin{proof}
Twierdzenie Taubesa zapewnia niezerowy niezmiennik,
a poprzednie twierdzenie o dodatniej krzywiźnie
wymuszałoby zerowanie wszystkich niezmienników.
Otrzymujemy sprzeczność.
\end{proof}

\section{Wymierna kontrakcja a niezmienniki gładkie}
\label{sec:tg-sw-rb}

\begin{twierdzenie}[Fintushel--Stern; formuła podniesienia, bez dowodu]
\label{thm:tg-sw-rbprinciple}
Niech $X$ będzie zamkniętą, po prostu spójną,
gładką, zorientowaną czterorozmaitością z
$b_2^+(X)>1$, zawierającą standardową konfigurację
$C_p$, $p\geq2$. Niech $W=X\setminus\operatorname{int}C_p$
i $X_{(p)}=W\cup B_p$ będzie wymierną kontrakcją
wzdłuż zanurzenia zachowującego orientację.
W tym sformułowaniu zakładamy dodatkowo
prostą spójność $X_{(p)}$, aby klasy wyznacznikowe
jednoznacznie określały struktury $\mathrm{Spin}^c$
na obu zamkniętych rozmaitościach.
Z twierdzenia \ref{thm:tg-rb-invariants} wynika
$b_2^+(X_{(p)})=b_2^+(X)>1$.
Ustalmy zgodne orientacje homologiczne.

Niech $\mathfrak s'$ będzie strukturą $\mathrm{Spin}^c$
na $X_{(p)}$ z $d_{X_{(p)}}(\mathfrak s')=0$.
Przy ustalonym sklejeniu brzegu nazywamy
\emph{dopuszczalnym podniesieniem} $\mathfrak s'$
strukturę $\mathfrak s$ na $X$, dla której
$\mathfrak s|_W\cong\mathfrak s'|_W$, restrykcje na
wspólnym brzegu zgadzają się jako struktury
$\mathrm{Spin}^c$, a ich przedłużenie przez $C_p$
spełnia $d_X(\mathfrak s)=d_{X_{(p)}}(\mathfrak s')=0$.
Ostatni warunek można równoważnie zapisać jako
$c_1(\mathfrak s)^2=c_1(\mathfrak s')^2-(p-1)$.
Każda taka struktura $\mathfrak s'$ ma
dopuszczalne podniesienie. Dla każdego
dopuszczalnego podniesienia zachodzi
\[
SW_{X_{(p)}}(\mathfrak s')=SW_X(\mathfrak s)
\]
przy odpowiednim zgodnym wyborze orientacji homologicznych.
Istnienie podniesień i równość niezmienników
przyjmujemy z pracy R.~Fintushela i R.~J. Sterna,
\emph{Rational blowdowns of smooth 4-manifolds},
J. Differential Geom. 46 (1997), 181--235,
propozycja 8.4 i twierdzenie 8.5
(wersja arXiv:alg-geom/9505018).
Na obu zamkniętych rozmaitościach $H^2(-;\mathbb Z)$
nie ma torsji. Sprawdźmy również zgodność samych
struktur na $W$. Niech $\gamma$ będzie generatorem
$\pi_1(\partial C_p)\cong\mathbb Z/p^2\mathbb Z$.
Z prostej spójności $X$ i twierdzenia van Kampena
wynika, że obraz $\gamma$ normalnie generuje
$\pi_1(W)$, gdyż $C_p$ jest po prostu spójna.
Zatem $H_1(W;\mathbb Z)$ jest cykliczna rzędu
$d$ dzielącego $p^2$. Z rachunku sklejenia
w stwierdzeniu \ref{prop:tg-rb-pi1-general}
i surjekcji $\pi_1(\partial B_p)\to\pi_1(B_p)$
wynika, że $H_1(X_{(p)};\mathbb Z)$ jest ilorazem
tej grupy przez podgrupę generowaną przez
$p$-krotność jej generatora. Ma więc rząd
$\gcd(d,p)$. Prosta spójność $X_{(p)}$ wymusza
$\gcd(d,p)=1$, a ponieważ $d\mid p^2$,
otrzymujemy $d=1$. Uniwersalne twierdzenie
o współczynnikach daje brak torsji również
w $H^2(W;\mathbb Z)$. Zatem zgodność wyznaczników
na $W$ zapewnia zgodność samych struktur,
co uzasadnia przejście od sformułowania źródłowego
do sformułowania w języku struktur $\mathrm{Spin}^c$.
Nie dowodzimy tutaj analitycznej formuły sklejenia.
\end{twierdzenie}

\begin{uwaga}[Dlaczego sformułowanie jest ostrożne]
Pełna formuła sklejenia wymaga zidentyfikowania struktur
$\mathrm{Spin}^c$ na brzegu oraz dopasowania ich przedłużeń
przez $B_{p,q}$, a także uwzględnienia wymiarów
oczekiwanych i orientacji przestrzeni moduli. Samo
porównanie $c^2$ lub wyznacznika formy przecięcia
nie wyznacza liczby rozwiązań równań cechowania.
Dla $b_2^+=1$ trzeba dodatkowo kontrolować komory,
których definicję i formułę zmiany przedstawimy w rozdziale
\ref{ch:tg-sw-analiza}.
\end{uwaga}

\begin{przyklad}[Kontrola równości wymiarów oczekiwanych]
Niech $X\to X_{p,q}$ będzie uogólnioną wymierną kontrakcją
konfiguracji długości $r$. Oznaczmy
$\kappa(Z):=2e(Z)+3\sigma(Z)$, a przez
$c\in H^2(X;\mathbb Z)$ i $c'\in H^2(X_{p,q};\mathbb Z)$
klasy wyznacznikowe struktur $\mathrm{Spin}^c$
$\mathfrak s$ i $\mathfrak s'$ na odpowiednich
rozmaitościach. Jeżeli
$\langle c'^2,[X_{p,q}]\rangle
=\langle c^2,[X]\rangle+r$, to ze wzoru indeksowego i
$\kappa(X_{p,q})=\kappa(X)+r$ wynika
$d_{X_{p,q}}(\mathfrak s')=d_X(\mathfrak s)$.
Użyliśmy wzoru z twierdzenia
\ref{thm:tg-general-numeric}; rachunek indeksu
nie wymaga założenia, że obie rozmaitości
są po prostu spójne.
Warunek o kwadratach klas nie jest automatyczny;
wybór dopuszczalnego podniesienia musi go zapewnić.
Ten rachunek wyjaśnia, jakie dane poza samą formą
przecięcia są potrzebne do porównania niezmienników.
\end{przyklad}

\begin{lemat}[Warunek wymiarowy dla wymiernej kontrakcji sfery o kwadracie $-4$]
\label{lem:tg-sw-c2-admissible-pairing}
Niech $X$ będzie zamkniętą, spójną, zorientowaną, gładką
czterorozmaitością zawierającą gładko zanurzoną sferę
$S$ o samoprzecięciu $S^2=-4$. Niech $C_2$ będzie jej
domkniętym otoczeniem dyskowym, $W=X\setminus
\operatorname{int}C_2$, a $X_{(2)}=W\cup_{\varphi}B_2$
wymierną kontrakcją. Załóżmy, że struktury
$\mathrm{Spin}^c$ $\mathfrak s$ na $X$ i $\mathfrak s'$
na $X_{(2)}$ mają zgodne ograniczenia do $W$.
Oznaczmy $c=c_1(\mathfrak s)$ i
$c'=c_1(\mathfrak s')$ oraz $m=\langle c,[S]\rangle$.
Wtedy
\[
\langle c^2,[X]\rangle-
\langle (c')^2,[X_{(2)}]\rangle=-\frac{m^2}{4}.
\]
W szczególności $d_X(\mathfrak s)=d_{X_{(2)}}(\mathfrak s')$
wtedy i tylko wtedy, gdy $|m|=2$.
Warunek ten jest konieczny dla dopuszczalnego
podniesienia z Twierdzenia~\ref{thm:tg-sw-rbprinciple},
lecz nie dowodzi istnienia takiego podniesienia.
\end{lemat}

\begin{proof}
Brzeg $C_2$ jest wymierną sferą homologiczną, a
$B_2$ wymierną kulą homologii. Ciągi
Mayera--Vietorisa nad $\mathbb Q$ rozkładają zatem
$H_2(X;\mathbb Q)$ na ortogonalną sumę klas z $W$
i prostą rozpiętą przez $[S]$, a
$H_2(X_{(2)};\mathbb Q)$ utożsamiają z częścią
pochodzącą z $W$. Zgodność ograniczeń struktur
$\mathrm{Spin}^c$ do $W$ zapewnia zgodność
składników klas $c$ i $c'$ na tej części.
Składowa Poincarégo dualna do $c$, leżąca na
prostej rozpiętej przez $[S]$, jest równa
$-\frac m4[S]$, gdyż $[S]^2=-4$ i jej
parowanie z $[S]$ ma wynosić $m$.
Jej kwadrat to $(-m/4)^2(-4)=-m^2/4$;
różnica kwadratów klas wynosi więc tyle samo.

Z rachunku niezmienników wymiernej kontrakcji mamy
$2e(X_{(2)})+3\sigma(X_{(2)})
=2e(X)+3\sigma(X)+1$.
Stąd, korzystając ze wzoru indeksowego,
\[
4\bigl(d_X(\mathfrak s)-d_{X_{(2)}}(\mathfrak s')\bigr)
=-\frac{m^2}{4}+1.
\]
Równość wymiarów zachodzi dokładnie dla $m^2=4$,
czyli $m=\pm2$. Dla klas charakterystycznych
$m$ jest parzyste, ponieważ $[S]^2=-4$.
Co więcej, jeżeli obie struktury $\mathrm{Spin}^c$
o zgodnych ograniczeniach rzeczywiście istnieją,
to obydwa wymiary oczekiwane są całkowite.
Ze wzoru
$d_X(\mathfrak s)-d_{X_{(2)}}(\mathfrak s')
=(4-m^2)/16$
wynika wtedy $m\equiv2\pmod4$.
Istnienie zgodnego przedłużenia jest więc
silniejszym warunkiem niż sama parzystość $m$;
w szczególności $m=0$ wyklucza takie przedłużenie.
Natomiast $m\equiv2\pmod4$ nie zapewnia jeszcze
równości wymiarów, która wymaga $m=\pm2$.
\end{proof}

\begin{przyklad}[Wymierna kontrakcja sekcji w $E(4)$]
Niech $E(4)$ oznacza standardową względnie minimalną,
po prostu spójną powierzchnię eliptyczną o
$e(E(4))=48$ i $\sigma(E(4))=-32$.
Ma ona $b_2^+=7$ i $b_2^-=39$.
Jej sekcja jest sferą o samoprzecięciu $-4$,
a więc stanowi konfigurację $C_2$. Po wymiernej
kontrakcji otrzymujemy gładką czterorozmaitość $Y$.
Jej prostą spójność przyjmujemy z przykładu
w pracy R.~Fintushela i R.~J. Sterna,
\emph{Six lectures on four 4-manifolds},
wykład 4, sekcja 3, s.~30.
Z twierdzenia \ref{thm:tg-rb-invariants} ma ona liczby $e(Y)=47$, $\sigma(Y)=-31$,
$b_2^+(Y)=7$ i $b_2^-(Y)=38$.

Powierzchnia $E(4)$ ma typ prosty Seiberga--Wittena,
a jej wielomian, z dokładnością do ogólnego znaku,
wynosi $(t-t^{-1})^2=t^2-2+t^{-2}$.
Korzystamy z obliczenia w pracy Fintushela--Sterna
\emph{Rational blowdowns of smooth 4-manifolds},
wzór (22) w wersji arXiv.
Tutaj $F$ oznacza zarówno klasę włókna, jak i
jej klasę Poincarégo dualną w kohomologii
(gdy piszemy o klasach wyznacznikowych),
a $t$ jest symbolem tej ostatniej klasy
w pierścieniu grupowym.
Klasy podstawowe to zatem $-2F,0,2F$,
o współczynnikach odpowiednio $1,-2,1$.
Typ prosty oznacza, że wszystkie niezerowe
niezmienniki występują dla struktur o wymiarze
oczekiwanym zero.
Gdy $S$ jest sekcją, zachodzi $F\cdot S=1$,
a zatem $(\pm2F)\cdot S=\pm2$.
Na mocy Lematu~\ref{lem:tg-sw-c2-admissible-pairing}
jedynymi spośród trzech wyjściowych klas podstawowych,
które spełniają konieczny warunek równości wymiarów
oczekiwanych, są $\pm2F$:
dla nich $\langle\pm2F,[S]\rangle=\pm2$.
Klasa zerowa ma parowanie $0$ i nie może być
dopuszczalnym podniesieniem w tym sklejeniu.
Aby przejść od konieczności tego warunku
do kompletnego obliczenia, korzystamy dodatkowo
z zachowania typu prostego przy wymiernej
kontrakcji (wniosek 8.6 tej samej pracy)
oraz z twierdzenia o zstępowaniu klas
w konfiguracji taut, w przytoczonym
wykładzie 4, sekcji 3, s.~30.
W przypadku $C_2=\{S\}$ warunek taut oznacza
$|k\cdot S|\leq2$ dla każdej klasy podstawowej
$k$ rozmaitości wyjściowej. Zachodzi on
dla wszystkich trzech wymienionych klas.
Twierdzenie zapewnia, że dokładnie klasy
o $|k\cdot S|=2$ zstępują jako klasy podstawowe
oraz że otrzymujemy wszystkie klasy podstawowe $Y$.
Zstępują więc $\pm2F$ z niezmiennikami równymi
$1$ przy zgodnym wyborze orientacji.
Oznaczając klasę zstępującą z $2F$ przez $k_Y$,
otrzymujemy $k_Y^2=1$ na mocy lematu
\ref{lem:tg-sw-c2-admissible-pairing}
i $F^2=0$. W szczególności $k_Y\ne-k_Y$.
W pierścieniu grupowym, dla symbolu $t_Y$ klasy
$k_Y$, otrzymujemy
$SW_Y=t_Y+t_Y^{-1}$ z dokładnością do ogólnego znaku.
Kompletność tego obliczenia jest konsekwencją
formuły sklejenia i twierdzenia o konfiguracji taut,
nie samego rachunku wymiarów.
\end{przyklad}

\section{Jak wykorzystać teorię do rozróżniania struktur gładkich?}
\label{sec:tg-sw-exotic}

\begin{stwierdzenie}[Zasada rozróżniania]
Niech $X_1,X_2$ będą zamkniętymi, spójnymi,
zorientowanymi, gładkimi czterorozmaitościami.
Załóżmy, że są homeomorficzne z zachowaniem orientacji
i że $b_2^+(X_i)>1$ dla $i=1,2$.
Jeżeli ich niezmienniki Seiberga--Wittena
w wymiarze oczekiwanym zero nie odpowiadają sobie
przy żadnym możliwym przeniesieniu struktur
$\mathrm{Spin}^c$ przez dyfeomorfizm zachowujący
orientację, po uwzględnieniu orientacji homologicznej
(i ewentualnej wspólnej zmiany znaku), to $X_1$
i $X_2$ nie są dyfeomorficzne z zachowaniem orientacji.
\end{stwierdzenie}

\begin{proof}
Gdyby istniał dyfeomorfizm zachowujący orientację,
indukowałby izometrię form przecięcia
na kohomologii całkowitej modulo torsja oraz
bijekcję struktur $\mathrm{Spin}^c$.
Niezmienniczość Seiberga--Wittena wymuszałaby
wówczas zgodność wartości niezmienników
po przeniesieniu struktur i orientacji homologicznej,
wbrew założeniu.
Jeżeli $H^2(-;\mathbb Z)$ jest wolna od torsji,
możemy do tego celu porównywać klasy wyznacznikowe
z ich niezmiennikami przy wszystkich izometriach
odpowiednich krat (jest to warunek wystarczający).
W obecności torsji sama izometria kraty
pomija część informacji o klasach wyznacznikowych,
a przy 2-torsji nawet pełne klasy wyznacznikowe
nie rozróżniają wszystkich struktur $\mathrm{Spin}^c$.
Dla $b_2^+=1$ potrzebna jest wersja komorowa
tej zasady, z przeniesieniem także komory.
\end{proof}

\begin{center}
\begin{tikzpicture}[>=Stealth,node distance=8mm]
  \node[draw=tgBlue,fill=tgBlue!9,rounded corners=3pt,
    minimum width=29mm,minimum height=11mm,align=center,font=\small] (a) at (-3,0)
    {graf plumbingowy\\$Q_{p,q}$};
  \node[draw=tgTeal,fill=tgTeal!12,rounded corners=3pt,
    minimum width=31mm,minimum height=11mm,align=center,font=\small] (b) at (.4,0)
    {wymierna kontrakcja\\$X\longmapsto X_{p,q}$};
  \node[draw=tgCoral,fill=tgCoral!12,rounded corners=3pt,
    minimum width=28mm,minimum height=11mm,align=center,font=\small] (c) at (3.9,0)
    {równania SW\\klasy podstawowe};
  \draw[tgTree,-{Stealth[length=2.7mm]}] (a)--(b);
  \draw[tgAccent,-{Stealth[length=2.7mm]}] (b)--(c);
  \node[font=\footnotesize,text=tgNavy,align=center] at (.4,-1.08)
    {Topologia: $\pi_1,Q,e,\sigma$\quad i osobno\quad gładkość: $SW$};
\end{tikzpicture}
\end{center}

\begin{uwaga}[Granice obecnego rozdziału]
Nie przeprowadziliśmy pełnego dowodu twierdzeń o zwartości
i transwersalności przestrzeni moduli, indeksu Atiyaha--Singera,
twierdzenia Taubesa ani formuły Fintushela--Sterna.
Przedstawiliśmy natomiast precyzyjnie niezbędne obiekty,
pełne rachunki topologiczne i logiczne miejsce, w którym
niezmienniki analityczne rozstrzygają problem gładkości.
W szczególności w przypadkach $b_2^+=1$, takich jak egzotyczne
$\mathbb{CP}^2\#7\overline{\mathbb{CP}}^{\,2}$, potrzebna
jest jeszcze osobna teoria komór i bardziej szczegółowa
formuła sklejenia. Tych kroków nie zastępujemy
rachunkiem charakterystyki Eulera.
\end{uwaga}


% Iteracja 11. Bez preambuły; wykorzystuje istniejące środowiska i barwy.
\begin{uwaga}[Wymiar oczekiwany a niezerowość niezmiennika]
Wartość $d(\mathfrak s)=0$ oznacza jedynie zerowy indeks
formalnego operatora deformacyjnego. Nie zapewnia istnienia
żadnego rozwiązania równań Seiberga--Wittena, regularności
wszystkich rozwiązań ani niezerowego zliczenia.
Do obliczenia niezmiennika potrzebne są niezależnie
zwartość, transwersalność i orientacje; przy $b_2^+=1$
również wybór komory. Dlatego sama zgodność klas
charakterystycznych i kwadratów ich pierwszych klas
Cherna nie stanowi dowodu równości niezmienników.
\end{uwaga}

